\chapter{Conclusion}
\label{ch:conclusion}

This thesis set out to realize the promise of energy-based controllable text generation, in which a frozen autoregressive language model is treated as an energy function and steered at inference time by sampling from it with gradient-guided Markov chain Monte Carlo. The appeal of that promise is real: a single frozen model, augmented with an additive constraint at generation time, could in principle be steered toward any objective without the cost of retraining. Across a systematic study of one hundred and forty-five sampling configurations, five energy functions, and two samplers, complemented by four diagnostic experiments, the thesis found that the central premise on which the promise rests does not hold. The frozen likelihood of an autoregressive language model does not admit local gradient-based navigation on discrete text: its gradient is not a usable search direction. The narrower phrasing is deliberate and is earned by the evidence, because the energy itself proves searchable by gradient-free means, a single top-$k$ rescoring pass and a Gibbs sampler both recovering coherent text on the same benchmark; what fails is not the evaluation of the energy but the attempt to follow its gradient.

The contribution of the thesis is not this negative observation alone, which the difficulty the field has had with these methods already hints at, but the mechanistic account that accompanies it, established by independent means and supported at every step by direct measurement. Following the true gradient was shown to be statistically indistinguishable from following a random direction of equal magnitude, on all five energy functions, once a normalization confound was removed by an ablation designed for that purpose. This gradient fallacy was explained by a linearization failure: the first-order surrogate that the discrete sampler uses to rank candidate tokens is uncorrelated with the true change in energy, because the smallest admissible discrete step is far larger than the radius over which the local linear model is valid, and the input-embedding gradient is in addition structurally blind to a substantial part of the energy change. The Metropolis--Hastings correction, applied faithfully rather than omitted, was shown to paralyze the continuous sampler, and the trace analysis located the cause precisely in the proposal term of the acceptance ratio, which collapses because the nearest-neighbour projection makes the drift discontinuous at every cell boundary. A likelihood trap was measured on the study's own models, showing that the lowest-energy text is the most degenerate, so that success at the optimization objective is failure at the underlying goal. Turning to amortized inference, the GFlowNet experiments exhibited three distinct failure modes, and a unifying experiment showed that Langevin sampling on the GFlowNet-tuned energy is indistinguishable from sampling on the untuned base, demonstrating that the pathology survives amortization and is therefore a property of the energy rather than of any sampler.
% PRIOR WORDING (Phase 3): A constrained-generation extension confirmed that the constraint gradient carries no reliable steering signal on this landscape.
A constrained-generation extension refined rather than merely repeated this picture: unlike the fluency gradient, the constraint gradient does carry a directional signal, but it does not produce reliable control on the discrete sampler, and where a signal appears it is entangled with the departure from the fluent manifold that a constraint-only objective causes.

Taken together, these findings locate the cause of the failure in the training objective of autoregressive language models, which shapes accurate local conditional distributions but never a smooth global energy over sequences. This diagnosis is not an endpoint but a signpost, because it makes a falsifiable prediction: a diffusion language model, whose objective does shape a score function over the data, should not fail in the same way.
% PRIOR WORDING (Phase 3, before the control was run): ...and the thesis sets out that positive-control experiment as the natural and immediate next step...
The thesis runs that positive control rather than only proposing it. A score-entropy discrete diffusion model on the same tokenizer, tested at two scales, supplies a clearly positive local direction where the autoregressive gradient was statistically zero, natively recovers the flipped token the gradient-guided sampler never recovered, and, in the sharpest form of the test, lifts recovery from zero to thirty-nine percent when it replaces only the proposal inside the thesis's own exact-energy chain; a trained noisy-state classifier steers an independent judge, with a reach bounded and setting-contingent for the same off-manifold reasons the constraint experiment exposed. These are pilots on smaller, differently trained models, and are read as such, but every arm of the prediction is borne out, which isolates the training objective as the causal variable and points to the complementary route of aligning a model's sequence density with a target energy by score matching. What the thesis provides to the larger effort of building inference-time controllable generation on frozen models is therefore a precise characterization of where and why the most common approach breaks down, drawn carefully enough to guide what comes next, together with a positive control that confirms the explanation. A negative result about a method, when it is accompanied by a mechanism, a measurement, and a falsifiable hypothesis, is a positive result about understanding, and it is in that spirit that this work is offered.
